% SOLO P2
% Texto de Marcia Arce (secciones 4.1 a 4.4 del documento compartido del grupo),
% pasado a LaTeX sin cambios de contenido: solo se adaptaron el formato de los
% decimales (punto, como en el resto del informe), las citas y los nombres de
% variables con \var{}.
% Sugerencia para la autora: en el último párrafo de la subsección D podría decir
% "a los conjuntos de validación y prueba", porque en P05 el mismo preprocesador
% ajustado con entrenamiento también transforma el conjunto de prueba.
\section{Preprocesamiento e ingeniería de características}
\label{sec:prepro}

\subsection{Tratamiento de valores faltantes}
La variable \var{TotalCharges}, originalmente almacenada como texto, se convirtió a formato numérico
mediante \var{pd.to\_numeric(errors=\textquotedbl{}coerce\textquotedbl{})}, identificándose 11 valores faltantes, equivalentes al
0.16\% de las observaciones. Se verificó que todos estos registros correspondían a clientes con una
antigüedad (\var{tenure}) de cero meses.

Para su tratamiento, se utilizó \var{SimpleImputer} con una estrategia de imputación constante de cero,
aplicada exclusivamente a \var{TotalCharges}~\cite{simpleimputer}. Esta decisión se fundamentó en el
supuesto de que los clientes sin meses completos de antigüedad podrían no haber acumulado cargos. Se
consideró más coherente con el significado de la variable que sustituir los valores faltantes por la
mediana del conjunto de entrenamiento (1385.30), lo que habría introducido cargos acumulados
considerables en estos registros. Sin embargo, la imputación adoptada constituye una hipótesis
metodológica, dado que los valores reales ausentes no pueden verificarse.

Las demás características numéricas y categóricas no presentaron valores faltantes, por lo que no se
aplicaron procedimientos adicionales de imputación.

\subsection{Codificación y escalamiento de características}
Las 16 características categóricas, incluida \var{SeniorCitizen}, se transformaron mediante
\emph{One-Hot Encoding}. Este procedimiento representa cada categoría mediante indicadores binarios,
evitando introducir relaciones ordinales artificiales entre categorías nominales~\cite{onehot}. Se
configuró \var{handle\_unknown=\textquotedbl{}ignore\textquotedbl{}} para permitir la transformación de categorías no identificadas
durante el ajuste.

Como resultado, las 19 características predictoras originales quedaron representadas mediante 46
columnas numéricas: tres correspondientes a variables cuantitativas y 43 generadas mediante
codificación categórica. Este incremento corresponde a un cambio en la representación de los datos y
no a la incorporación de información adicional sobre los clientes.

Para la regresión logística se incorporó \var{StandardScaler} sobre las tres características numéricas.
Esta técnica estandariza las variables utilizando la media y desviación estándar del conjunto de
entrenamiento, reduciendo las diferencias de escala que pueden afectar la optimización y la
regularización de modelos lineales~\cite{standardscaler}. En contraste, para el árbol de decisión se
conservaron las escalas originales, dado que sus particiones se basan en umbrales de las
características y no requieren estandarización.

\subsection{Análisis y decisiones de ingeniería de características}
Se examinó la posible redundancia entre las características cuantitativas mediante correlaciones de
Pearson calculadas exclusivamente sobre el conjunto de entrenamiento, aplicando previamente la
imputación de cero definida para \var{TotalCharges}.

Se obtuvieron correlaciones de 0.834 entre \var{tenure} y \var{TotalCharges}, y de 0.650 entre
\var{MonthlyCharges} y \var{TotalCharges}. Asimismo, la correlación entre \var{TotalCharges} y el
producto \var{tenure} $\times$ \var{MonthlyCharges} fue de 0.9996, evidenciando una asociación lineal
muy elevada entre estas cantidades, aunque no una equivalencia exacta.

Estos resultados sugieren una posible redundancia informativa. Sin embargo, la correlación entre
predictores no permite determinar por sí sola si alguna característica carece de utilidad para la
clasificación. La selección de características debe considerar su contribución al desempeño del
modelo y los riesgos asociados al sobreajuste~\cite{kuhn}.

Por ello, se decidió conservar inicialmente las 19 características originales, sin aplicar
eliminación de predictores basada únicamente en correlaciones. La ingeniería realizada se centró en
transformar su representación mediante codificación y escalamiento, sin incorporar nuevas
características derivadas. La conveniencia de eventuales modificaciones deberá determinarse mediante
validación experimental.

\subsection{Implementación del preprocesamiento}
Las transformaciones se organizaron mediante \var{Pipeline} y \var{ColumnTransformer} de scikit-learn,
implementándose la función \var{make\_preprocessor()} para construir dos preprocesadores reutilizables,
correspondientes a regresión logística y árbol de decisión. Ambos comparten la imputación constante
de cero para \var{TotalCharges} y la codificación categórica, diferenciándose en la aplicación de
\var{StandardScaler} únicamente para la regresión logística.

Para prevenir la fuga de información (\emph{data leakage}), las transformaciones que requieren
aprender parámetros de los datos, como la identificación de categorías y el cálculo de estadísticas
de estandarización, se ajustaron exclusivamente sobre el conjunto de entrenamiento.
Las transformaciones aprendidas se aplicaron posteriormente a validación y prueba sin reajustarlas,
siguiendo las recomendaciones de scikit-learn~\cite{leakage}.

